% !TeX program = xelatex
% 题源根目录：references/302模拟27md/
% 武忠祥基础：2027模考/武忠祥/27考研基础摸底测试（高数、线代）/27考研基础摸底测试（高数、线代）数二.pdf.md
% 武忠祥8月：2027模考/武忠祥/武老师8月突击强化测试/8月强化突击测试.pdf.md
% 李林：2027模考/李林《六套卷》/数学二/试卷/考研数学冲刺预测6套卷.pdf.md
% 张宇：2027模考/张宇《八套卷》/张宇八套卷/数二/【紧凑】张宇八套卷-数二（已勘误）.pdf.md
% 李艳芳：2027模考/李艳芳《三套卷》/数学二/【紧凑版】李艳芳三套卷做题本 .pdf.md
\documentclass[UTF8,zihao=-4,fontset=fandol]{ctexart}
\usepackage[paperwidth=800bp,paperheight=600bp,left=48bp,right=48bp,top=40bp,bottom=40bp]{geometry}
\usepackage{amsmath,amssymb,xcolor,fancyhdr,enumitem}
\usepackage[hidelinks,bookmarksopen=true]{hyperref}
\renewcommand{\lim}{\operatorname*{lim}\limits}
\everymath{\displaystyle}
\definecolor{accent}{HTML}{175C57}
\setlength{\parindent}{0pt}\setlength{\parskip}{5bp}\setlength{\headheight}{15bp}
\setlength{\abovedisplayskip}{6bp}\setlength{\belowdisplayskip}{6bp}
\pagestyle{fancy}\fancyhf{}\renewcommand{\headrulewidth}{0pt}
\fancyfoot[L]{\footnotesize\textcolor{gray}{中值定理阶段模拟卷 3 · 4:3 横屏}}
\fancyfoot[R]{\footnotesize\thepage}\raggedbottom
\newcommand{\Question}[4]{\clearpage\def\CurrentNumber{#1}\def\CurrentSource{#4}\pdfbookmark[1]{第#1题}{q#1}{\small\color{accent}模拟卷 3\hfill #2 · #3 分}\par{\Large\bfseries 第 #1 题}\par\medskip\fontsize{14}{22}\selectfont}
\newcommand{\Solution}{\clearpage\fontsize{12}{18}\selectfont{\large\bfseries 第 \CurrentNumber\ 题解析}\par{\footnotesize\color{gray}题源：\CurrentSource}\par\smallskip}
\newcommand{\Known}{\par\textbf{已知与目标}\quad}
\newcommand{\Method}{\par\textbf{方法与条件}\quad}
\newcommand{\Derive}{\par\textbf{推导}\quad}
\renewcommand{\Check}{\par\textbf{结论与检查}\quad}
\newcommand{\Choices}[4]{\par\medskip\begin{minipage}[t]{.48\linewidth}(A)\enspace #1\end{minipage}\hfill\begin{minipage}[t]{.48\linewidth}(B)\enspace #2\end{minipage}\par\medskip\begin{minipage}[t]{.48\linewidth}(C)\enspace #3\end{minipage}\hfill\begin{minipage}[t]{.48\linewidth}(D)\enspace #4\end{minipage}\par}
\newcommand{\blank}{\underline{\hspace{6em}}}
\begin{document}
\pdfbookmark[0]{模拟卷3}{cover}
{\small\color{accent}工科数学分析基础 · 第三版上册}\par\vspace{20bp}
{\Huge\bfseries 阶段模拟卷 3}\par\medskip
{\Large 综合极限、函数关系与数列}\par\vspace{18bp}
\textbf{满分 100 分\qquad 建议用时 120 分钟}

考试范围：第一章；第二章第一至第三节，以及第四节中的 4.1 极值必要条件、4.2 Rolle、Lagrange、Cauchy 中值定理。教材对应 PDF 第 24–153 页。解答只使用这一范围内的方法。

第 1–4 题为单项选择，每题 5 分；第 5–8 题为填空，每题 5 分；第 9–12 题为解答，每题 15 分。解答题请写出理由。

每道题单独一页，下方空白供手写；下一页为该题解析。限时作答时请跳过解析页。题目总体为中等偏上，证明题和条件辨析题具有一定挑战。

题目选自本地 references/302模拟27md，具体作者、卷次与原题号见各题解析。改题型、增设任务或只取原题小问均逐题标注。

\Question{1}{单项选择}{5}{李林六套卷数二，第五套第 2 题；原题}
设 $f(x)$ 为连续函数，且 $\lim_{x\to1}\frac{f(2x-1)}{x-1}=2$，则
\[
\lim_{x\to0}\frac{f[(1+\sin x)^2]-f(1+\sin x)}x=\text{（\quad）}.
\]
\Choices{$0$}{$2$}{$1$}{$-1$}
\Solution
\Known 已知的是在 $1$ 附近的复合差商，要求另两个复合函数之差的极限。
\Method 先换元求 $f(1)$ 与 $f'(1)$，再分别以 $1$ 为基准拆出导数差商。
\Derive 令 $u=2x-1$，则 $x-1=\frac{u-1}{2}$。给定条件成为
\[
\lim_{u\to1}\frac{f(u)}{u-1}=1.
\]
连续性与有限极限给出 $f(1)=0$，进而 $f'(1)=1$。记 $s=\sin x$，则
\[
\begin{aligned}
\frac{f((1+s)^2)-f(1+s)}x
&=\frac{f((1+s)^2)-f(1)}{(1+s)^2-1}\frac{2s+s^2}x\\
&\quad-\frac{f(1+s)-f(1)}s\frac sx.
\end{aligned}
\]
当 $x\to0$，两个差商均趋于 $1$，其余两个比值分别趋于 $2$ 和 $1$，故所求极限为 $1$。
\Check 选 \textbf{C}。取 $f(u)=u-1$ 代入，原式直接成为 $\frac{\sin x+\sin^2x}x\to1$，可核对换元中的系数。

\Question{2}{单项选择}{5}{武忠祥8月强化突击测试，第 2 题；原题}
设 $f(x)$ 在 $x=a$ 处可导，则 $|f(x)|$ 在 $x=a$ 处不可导的充分条件是（\quad）。
\Choices{$f(a)=0$，且 $f'(a)=0$。}{$f(a)=0$，且 $f'(a)\ne0$。}{$f(a)>0$，且 $f'(a)>0$。}{$f(a)<0$，且 $f'(a)<0$。}
\Solution
\Known 已知 $f$ 自身可导，要判断取绝对值后何时出现不可导。
\Method $f(a)\ne0$ 时连续性使附近符号固定，因此只会保留 $f$ 或变为 $-f$。真正需要检查的是 $f(a)=0$ 时的左右差商。
\Derive 若 $f(a)=0$，则 $\frac{f(a+h)}h\to f'(a)$。由于
\[
\frac{|f(a+h)|-|f(a)|}h
=\begin{cases}\left|\frac{f(a+h)}h\right|,&h>0,\\-\left|\frac{f(a+h)}h\right|,&h<0,\end{cases}
\]
右导数为 $|f'(a)|$，左导数为 $-|f'(a)|$。$f'(a)\ne0$ 时两者不同，故 $|f|$ 不可导。

若 $f'(a)=0$，两侧差商反而都趋于 $0$，此时 $|f|$ 可导。$f(a)>0$ 或 $f(a)<0$ 的情况已由符号固定排除。
\Check 选 \textbf{B}。以 $f(x)=x-a$ 检验，$|f(x)|=|x-a|$ 的左、右导数确为 $-1,1$。

\Question{3}{单项选择}{5}{李艳芳三套卷数二，卷二第 1 题；原题}
若 $\lim_{x\to0}\frac{\tan x+xf(x)}{x^3}=\frac23$，则当 $x\to0$ 时，与 $f(x)+1$ 等价的无穷小量为（\quad）。
\Choices{$\cos x-1$}{$\ln\left(1+\frac x3\right)$}{$1-\mathrm e^{x^2}$}{$\sqrt[3]{1+x^2}-1$}
\Solution
\Known 给定三次差商，要从中分离出 $f(x)+1$ 的主要小量。
\Method 写成 $\tan x-x+x(f(x)+1)$，并用 Cauchy 中值定理计算 $\tan x-x$ 的三次系数。
\Derive 对 $\tan t-t$ 与 $t^3$ 在 $0,x$ 之间用 Cauchy 定理，存在中间点 $u$，使
\[
\frac{\tan x-x}{x^3}=\frac{\tan^2u}{3u^2}\to\frac13.
\]
因此
\[
\frac{f(x)+1}{x^2}
=\frac{\tan x+xf(x)}{x^3}-\frac{\tan x-x}{x^3}\to\frac13.
\]
令 $v=\sqrt[3]{1+x^2}$，由 $v^3-1=x^2$ 可知
\[
\frac{\sqrt[3]{1+x^2}-1}{x^2}=\frac1{v^2+v+1}\to\frac13.
\]
两者与 $x^2$ 的比值同为非零数 $\frac13$，所以二者等价。
\Check 选 \textbf{D}。A 的二次系数为 $-\frac12$，C 为 $-1$，B 是一次小量，都与所得正的二次项不同。

\Question{4}{单项选择}{5}{张宇八套卷数二，卷八第 1 题；原题}
设 $f$ 在 $(-1,1)$ 内有定义且在 $0$ 连续。判断下列结论：

（1）若 $\lim_{x\to0}\frac{f(x)}{\sqrt[3]x}=0$，则 $f$ 在 $0$ 处可导。

（2）若 $\lim_{x\to0}\frac{f(x)}{x^2}=0$，则 $f$ 在 $0$ 处可导。

（3）若 $f$ 在 $0$ 处可导，则 $\lim_{x\to0}\frac{f(x)}{\sqrt[3]x}=0$。

所有正确结论的序号为（\quad）。
\Choices{（1）}{（2）}{（2）（3）}{（1）（2）}
\Solution
\Known 三个命题都涉及导数与小量的关系，需要分别验证，不能擅自添加 $f(0)=0$。
\Method 正确命题用导数定义证明；错误命题用满足前提但不满足结论的函数，即反例，进行排除。
\Derive （1）取 $f(x)=|x|^{\frac23}$。有 $\frac{f(x)}{\sqrt[3]x}\to0$，但 $\frac{f(x)-f(0)}x$ 的右侧趋于 $+\infty$、左侧趋于 $-\infty$，故不可导。结论（1）错误。

（2）给定条件推出 $f(x)\to0$，连续性给出 $f(0)=0$。于是
\[
\frac{f(x)-f(0)}x=x\frac{f(x)}{x^2}\to0,
\]
故 $f'(0)=0$。结论（2）正确。

（3）取 $f(x)\equiv1$，函数处处可导，但 $\frac1{\sqrt[3]x}$ 不趋于 $0$。结论（3）错误。
\Check 选 \textbf{B}。两个反例在整个 $(-1,1)$ 内都有定义且在 $0$ 连续；原题确实没有给出 $f(0)=0$，这一条件不能自行补入。

\Question{5}{填空}{5}{张宇八套卷数二，卷六第 11 题；原题}
\[
\lim_{x\to0}\left(\frac{2+x}{\sin x}-\frac2x\right)=\blank.
\]
\Solution
\Known 两个发散项相减，必须先合并或拆成有有限极限的部分。
\Method 分离出 $\frac{x}{\sin x}$，其余部分由中值定理估计 $x-\sin x$ 的大小。
\Derive 严格变形为
\[
\frac{2+x}{\sin x}-\frac2x
=\frac{x}{\sin x}+\frac{2(x-\sin x)}{x\sin x}.
\]
第一项趋于 $1$。对 $t-\sin t$ 用 Lagrange 定理，可得 $x-\sin x=x(1-\cos u)$，其中 $|u|\le|x|$。又由三角恒等式有
\[
0\le1-\cos u=2\sin^2\frac u2\le\frac{u^2}2\le\frac{x^2}2.
\]
所以 $\left|\frac{2(x-\sin x)}{x\sin x}\right|\le\frac{x^2}{|\sin x|}\to0$，因为 $\frac{|\sin x|}{|x|}\to1$。
\Check 填 $\boxed{1}$。计算覆盖 $x\to0^+$ 和 $x\to0^-$，没有将两个无穷大直接相减。

\Question{6}{填空}{5}{李林六套卷数二，第一套第 1 题；选择题改为填空}
当 $x\to0$ 时，$ax+bx^2+\ln(1+x)$ 与 $\mathrm e^{x^2}-\cos x$ 是等价无穷小，则
\[
a=\blank,\qquad b=\blank.
\]
\Solution
\Known 右式是二次小量，要求左式一次项抵消后，二次项系数与之相同。
\Method 先除以 $x$，再除以 $x^2$。用 Cauchy 中值定理求 $\ln(1+x)-x$ 的二次系数。
\Derive 右式可分成 $(\mathrm e^{x^2}-1)+(1-\cos x)$，故
\[
\frac{\mathrm e^{x^2}-\cos x}{x^2}\to1+\frac12=\frac32.
\]
左式除以 $x$ 的极限为 $a+1$，必须为零，得 $a=-1$。令 $h(x)=\ln(1+x)-x$，在 $0,x$ 之间用 Cauchy 定理，得中间点 $u$，满足
\[
\frac{h(x)}{x^2}=\frac{\frac1{1+u}-1}{2u}
=-\frac1{2(1+u)}\to-\frac12.
\]
所以左式除以 $x^2$ 趋于 $b-\frac12$。令它等于 $\frac32$，得 $b=2$。
\Check 填 $\boxed{-1}$、$\boxed{2}$。代回后两式都与 $\frac32x^2$ 等价，保证两者之比趋于 $1$。

\Question{7}{填空}{5}{武忠祥基础摸底测试数二，第 11 题；原题}
设 $y=y(x)$ 是由方程 $xy+\mathrm e^y=x+1$ 确定的隐函数，则
\[
\left.\frac{\mathrm d^2y}{\mathrm dx^2}\right|_{x=0}=\blank.
\]
\Solution
\Known $y$ 通过含 $x,y$ 的方程确定，要求它在 $0$ 处的二阶导数。
\Method 把 $y$ 视为 $x$ 的函数，对等式两边连续求导。先求点值和一阶导数，再代入二阶关系。
\Derive 令 $x=0$，有 $\mathrm e^{y(0)}=1$，故 $y(0)=0$。第一次求导，得
\[
y+xy'+\mathrm e^y y'=1,
\qquad (x+\mathrm e^y)y'=1-y.
\]
在 $x=0,y=0$ 处，系数 $x+\mathrm e^y=1\ne0$，所以 $y'(0)=1$。再次求导时，$\mathrm e^y y'$ 是乘积，故
\[
2y'+xy''+\mathrm e^y(y')^2+\mathrm e^y y''=0.
\]
代入点值与一阶导数，得到 $2+y''(0)+1=0$，因此 $y''(0)=-3$。
\Check 填 $\boxed{-3}$。已分别核对点值满足原方程、一阶关系满足 $1=1$、二阶关系满足 $2-3+1=0$。

\Question{8}{填空}{5}{张宇八套卷数二，卷一第 11 题；原题}
\[
\lim_{x\to0}\left(\frac{\mathrm e^x-1}{\sin x}\right)^{\cot x}=\blank.
\]
\Solution
\Known 底数趋于 $1$、指数无界。充分靠近 $0$ 时分子分母同号，底数为正，可以取对数。
\Method 先求对数的极限。用 Cauchy 中值定理得到底数偏离 $1$ 的一次系数，再取指数还原。
\Derive 令 $h(x)=\mathrm e^x-1-\sin x$。有 $h(0)=0$，Cauchy 定理给出中间点 $u$，使
\[
\frac{h(x)}{x^2}=\frac{\mathrm e^u-\cos u}{2u}
=\frac{\mathrm e^u-1}{2u}+\frac{1-\cos u}{2u}\to\frac12.
\]
令 $z(x)=\frac{\mathrm e^x-1}{\sin x}-1=\frac{h(x)}{\sin x}$，则
\[
\frac{z(x)}x=\frac{h(x)}{x^2}\frac{x}{\sin x}\to\frac12.
\]
原幂式的对数为
\[
\cot x\ln(1+z(x))
=\frac{\ln(1+z(x))}{z(x)}\frac{z(x)}x(x\cot x)\to1\cdot\frac12\cdot1=\frac12.
\]
这里用到 $\frac{\ln(1+z)}z\to1$，在 $z=0$ 时按连续延拓理解。
\Check 填 $\boxed{\mathrm e^{\frac12}}$。底数为正且左右两侧的对数极限相同，取指数后双侧极限成立。

\Question{9}{解答}{15}{武忠祥8月强化突击测试，第 3 题；原题，分值重设}
设函数 $f(x)$ 在 $x=0$ 处连续，且
\[
\lim_{x\to0}\frac{xf(x)-\mathrm e^{2\sin x}+1}{\ln(1+x)+\ln(1-x)}=-3.
\]
证明 $f$ 在 $0$ 处可导，并求 $f'(0)$。
\Solution
\Known 要求从给定极限证明导数存在，不能预先对 $f$ 求导。
\Method 分母合成 $\ln(1-x^2)$；先确定点值，再用 Cauchy 中值定理计算指数部分的二次余量。
\Derive 记分子为 $N(x)$。因 $\frac{\ln(1-x^2)}{x^2}\to-1$，得 $\frac{N(x)}{x^2}\to3$，所以 $\frac{N(x)}x\to0$。又
\[
\frac{\mathrm e^{2\sin x}-1}x
=\frac{\mathrm e^{2\sin x}-1}{2\sin x}\frac{2\sin x}x\to2.
\]
由连续性可知 $f(0)=2$。令 $r(x)=\mathrm e^{2\sin x}-1-2x$，对 $r(t)$ 与 $t^2$ 用 Cauchy 定理，得中间点 $u$，使
\[
\begin{aligned}
\frac{r(x)}{x^2}
&=\frac{2\cos u\,\mathrm e^{2\sin u}-2}{2u}\\
&=\frac{(\cos u-1)\mathrm e^{2\sin u}}u+\frac{\mathrm e^{2\sin u}-1}u\to0+2=2.
\end{aligned}
\]
最后，由 $N(x)=x[f(x)-2]-r(x)$ 得
\[
\frac{f(x)-f(0)}x=\frac{N(x)}{x^2}+\frac{r(x)}{x^2}\to3+2=5.
\]
\Check 已证明可导，且 $\boxed{f'(0)=5}$。将 $f(x)=2+5x$ 代入，分子与 $3x^2$ 等价、分母与 $-x^2$ 等价，商确实趋于 $-3$。

\Question{10}{解答}{15}{李林六套卷数二，第六套第 20 题第（I）问；只取求函数的小问}
设 $f$ 在 $(0,+\infty)$ 内有定义，$f'(1)=1$，且对任意正数 $x,y$，有
\[
f(xy)=\frac{f(x)}y+\frac{f(y)}x.
\]
求 $f(x)$。
\Solution
\Known 只有 $1$ 处的导数已知，要求整个正半轴上的函数表达式。
\Method 先将函数关系乘以 $xy$，转成乘积对应求和的形式；用 $1$ 处的导数推出任意点可导，再用 Lagrange 定理的推论“导数恒为零的函数是常数”。
\Derive 令 $x=y=1$，得到 $f(1)=2f(1)$，所以 $f(1)=0$。定义 $F(x)=xf(x)$，则
\[
F(xy)=F(x)+F(y),\qquad F(1)=0,\qquad F'(1)=f(1)+f'(1)=1.
\]
固定任意 $x>0$。当 $h\to0$ 时，令 $y=1+\frac hx>0$，则
\[
\frac{F(x+h)-F(x)}h
=\frac{F(1+\frac hx)-F(1)}{h}
=\frac1x\frac{F(1+\frac hx)-F(1)}{\frac hx}\to\frac1x.
\]
所以任意 $x>0$ 都有 $F'(x)=\frac1x$。函数 $G(x)=F(x)-\ln x$ 因而满足 $G'(x)=0$。在任意两个正数之间使用 Lagrange 定理，可知 $G$ 为常数。又 $G(1)=0$，故 $F(x)=\ln x$。
\Check $\boxed{f(x)=\frac{\ln x}{x},\ x>0}$。直接代入函数关系，两边均为 $\frac{\ln x+\ln y}{xy}$；且 $f'(1)=1$，全部条件满足。

\Question{11}{解答}{15}{李艳芳三套卷数二，卷二第 21 题第（1）问；只取极值点位置估计}
设 $f$ 在 $[0,1]$ 上三阶可导，$f(0)=f(1)$，并在 $x=c$（$0<c<1$）处取得极值。已知对任意 $x\in(0,1)$，
\[
1\le|f''(x)|\le4.
\]
证明：$\frac13\le c\le\frac23$。
\Solution
\Known 极值点在内部，故 Fermat 定理给出 $f'(c)=0$；两个端点函数值相同。
\Method 用 Rolle 定理把端点到极值点的函数增量写成二阶导数乘距离平方，再比较两侧距离。三阶可导保证 $f''$ 在内部连续；它不能在不经过零的情况下改变符号。
\Derive 因 $|f''|\ge1$，所以 $f''$ 符号固定。必要时把 $f$ 换成 $-f$，可设 $1\le f''\le4$。

固定 $t=0$ 或 $t=1$，令 $K=\frac{f(t)-f(c)}{(t-c)^2}$。函数 $H(s)=f(s)-f(c)-K(s-c)^2$ 满足 $H(c)=H(t)=0$、$H'(c)=0$。先对 $H$ 用 Rolle 定理得到 $H'(u)=0$，再对 $H'$ 在 $c,u$ 间用 Rolle 定理，得到 $f''(v)=2K$。

因此存在 $u\in(0,c)$、$v\in(c,1)$，满足
\[
f(0)-f(c)=\frac12f''(u)c^2,\qquad
f(1)-f(c)=\frac12f''(v)(1-c)^2.
\]
由两端点值相同，得到 $f''(u)c^2=f''(v)(1-c)^2$。利用 $1\le f''\le4$，有
\[
c^2\le4(1-c)^2,\qquad (1-c)^2\le4c^2.
\]
因 $c,1-c>0$，开平方得 $c\le2(1-c)$、$1-c\le2c$，即所需范围。
\Check $\boxed{\frac13\le c\le\frac23}$。取 $f(x)=(x-\frac12)^2$，有 $f''=2$、$c=\frac12$，满足所有条件且位于估计区间。

\Question{12}{解答}{15}{张宇八套卷数二，卷八第 5 题；选择题改为证明与求极限}
设数列 $\{a_n\}$ 满足 $a_1>2$，且
\[
a_{n+1}=a_n-\frac{a_n-3}{a_n-2},\qquad n=1,2,\ldots.
\]
证明该数列收敛，并求其极限。
\Solution
\Known 初值仅要求大于 $2$，递推式的分母可能带来问题，必须先证明各项有定义。
\Method 换元 $b_n=a_n-2$，先证明从第二项起有统一下界，再用单调有界定理：单调且有界的实数数列必收敛。
\Derive 当 $b_n>0$ 时，递推式化为
\[
b_{n+1}=b_n-1+\frac1{b_n}.
\]
由于 $b+\frac1b\ge2$（可由 $\frac{(b-1)^2}{b}\ge0$ 得出），故 $b_{n+1}\ge1$。从 $b_1>0$ 出发可归纳知所有项都有定义，且 $b_n\ge1$ 对 $n\ge2$ 成立。

当 $n\ge2$ 时，
\[
b_{n+1}-b_n=-1+\frac1{b_n}\le0.
\]
所以从第二项起数列单调不增且有下界 $1$，存在极限 $L\ge1$。递推式中的倒数在 $L\ge1$ 处连续，取极限得
\[
L=L-1+\frac1L,\qquad L=1.
\]
因此 $a_n=b_n+2\to3$。有限个初始项不影响数列极限。
\Check $\boxed{\lim_{n\to\infty}a_n=3}$。$a_1=3$ 时为常数列；即使 $2<a_1<3$，第二项也至少为 $3$，已覆盖全部允许初值。
\end{document}
